\section{Introduction}
\subsection{Motivation}
Biodiversity is defined as ``the variability among living organisms from all sources including, inter alia, terrestrial, marine and other aquatic ecosystems and the ecological complexes of which they are part; this includes diversity within species, between species and of ecosystems.''\footnote{Convention on Biological Diversity (1993) \textit{Article 2. Use of terms}. Available at: \url{https://www.cbd.int/convention/articles/?a=cbd-02} (Accessed: 8 April 2026).} In this sense, biodiversity refers to the biotic component of ecological systems: the variety and variability of living organisms, from plants and animals to fungi, bacteria, and protists. Biodiversity is integral to ecosystem functioning and stability, supporting processes such as nutrient cycling, pollutant breakdown, water regulation, carbon storage, and ecosystem productivity \parencite{Stork2009, cardinale2012, gascon2015}. Furthermore, biodiversity is fundamental to human survival and well-being. It supports food security through pollination, soil fertility, fisheries, and livestock systems \parencite{gascon2015,brauman2020}. It also provides biological materials and natural compounds used in medicine and drug discovery, including antibiotics such as penicillin \parencite{gascon2015}. Beyond these material contributions, biodiversity supports regulating ecosystem services such as water purification, climate regulation, and carbon storage, while also contributing to biomass-based energy systems and sustaining human health, livelihoods, and well-being \parencite{brauman2020, naeem2016, kim2024}.

Despite the overwhelming benefits of biodiversity to humanity and ecosystems, biodiversity is declining at an unprecedented rate. Approximately 74\% of the world's terrestrial wilderness has experienced degradation, and approximately 9.6\% has been completely lost since the early 1990s \parencite{watson2016}. Marine ecosystems have also been heavily affected, with 66\% of the ocean's surface having been significantly impacted by human activities \parencite{diaz2019}. Moreover, it is reasonably likely that most coral reefs will be lost if global temperatures exceed 2$^\circ$C \parencite{frieler2012}. Furthermore, human activities threaten the survival of approximately 1 million species, with approximately 25\% of plants and animals threatened \parencite{diaz2019}. Species extinction rates have also increased and are now estimated to be around 1,000 times higher than the natural background rate \parencite{pimm2014}. Compounding this challenge, conservation efforts are underfunded, with a deficit exceeding half a trillion US dollars annually \parencite{Deutz2020}. The acceleration of biodiversity loss and the drastic underfunding of conservation efforts have necessitated a shift towards the prioritization of biodiversity conservation and ecological rehabilitation.

Many initiatives around the world are being adopted to combat biodiversity loss and climate change, which are intricately linked. These initiatives focus on international cooperation, expanding protected areas, reducing pollution, and providing incentives for sustainable practices 
\parencite{IPBES2019spm, IPCC2023}. Specific policy initiatives such as the 30x30 target---which aims to protect 30\% of the world's land area by 2030---represent critical efforts for ecological recovery. If successful, they could generate large gains and rebounds in biodiversity. Achieving these goals would have significant impacts, including carbon sequestration, improved water quality regulation, and expanded capabilities to mitigate nutrient pollution \parencite{zeng2022}. For these initiatives to be successful, it is crucial to have accurate and robust models that can predict biodiversity patterns and trends. These models can inform conservation strategies, identify priority areas for protection, and assess the potential impacts of different conservation actions \parencite{sinclair2010}.

In this thesis, we use the broad term \emph{Spatial Biodiversity Models} to refer to models that predict one or more biodiversity measurements across space. This terminology emphasizes the spatially explicit nature of these models, which relate biodiversity patterns to environmental conditions at particular locations \parencite{ferrier2006}. Moreover, models that predict community-level biodiversity metrics are referred to as community-level biodiversity models or community-level models in the literature \parencite{ferrier2002, nieto2018}. These community-level biodiversity metrics can include species richness, total, or relative abundance, intactness, or compositional similarity \parencite{lamb2009, santini2017}. Broadly, community-level modeling approaches can be grouped into three main strategies based on how they produce predictions for community-level biodiversity metrics \parencite{ferrier2006}: \begin{inparaenum}[(1)]\item Bottom-up (``predict first, assemble later'') approaches, which commonly use species distribution models (SDMs) to predict the spatial distribution of individual species based on environmental variables; multiple SDMs can be aggregated or stacked to estimate community-level metrics \parencite{ferrier2006, distler2015, beery2021}; \item sideways or ``assemble and predict together'' approaches, which frequently use joint species distribution models (JSDMs) to model multiple species simultaneously, accounting for shared environmental responses and residual species associations \parencite{warton2015,wilkinson2019,pollock2020}; and \item top-down (``assemble first, predict later'') approaches, in which the model directly predicts community-level biodiversity metrics rather than first modeling individual species or assemblages \parencite{ferrier2006, pollock2020}.\end{inparaenum} 

While each strategy has its own strengths and weaknesses and excels in different contexts, it is unlikely that one modeling strategy will consistently outperform the others \parencite{ferrier2006, zhang2019}. However, bottom-up approaches are the most extensively researched modeling strategy for community-level biodiversity prediction. Many studies have focused on improving species distribution models, including benchmark evaluations \parencite{norberg2019, valavi2022}, reviews of specific models \parencite{waldock2022}, and guidance on best practices \parencite{araujo2019}. Similarly, ``assemble and predict together'' models have also been the focus of many studies. Joint species distribution models have been developed partly to address the limited treatment of species co-occurrence and biotic interactions in traditional species distribution models \parencite{wilkinson2019, escamilla2022}. Top-down models, on the other hand, are less widely studied and underutilized \parencite{pollock2020}. A significant gap remains in our understanding, comparison, and evaluation of bottom-up and top-down models, both between these modeling approaches and within the different methods they encompass. This highlights the need for further research and methodological development of top-down models \parencite{ferrier2006, pollock2020}.

Currently, top-down models are often evaluated in individual studies that use different datasets, model types, biodiversity metrics, and evaluation procedures \parencite{mondal2021, cai2023, ngor2023, baggstrom2025}. This can make it difficult to understand the current landscape of top-down models. While there have been some reviews of community-level modeling compared to other spatial biodiversity models \parencite{nieto2018}, there remains a lack of best practice guidelines for top-down models, comparable to those available for bottom-up approaches \parencite{araujo2019}. This leaves a methodological gap in the development, comparison, and rigorous evaluation of top-down modeling approaches \parencite{pollock2020}. A flexible modeling pipeline could help address this gap by enabling more systematic evaluation of different top-down models and by clarifying their strengths and limitations for predicting community-level biodiversity patterns and trends. Furthermore, a pipeline would enable a more nuanced evaluation of these models and their applicability across different ecological contexts, and allow for quick comparison across different models. This would ultimately contribute to the advancement of biodiversity modeling as a whole and a deeper exploration of the potential of top-down models.


\subsection{Research Objectives and Scope}\label{sec:research_objectives}

The primary objective of this Master's thesis is to develop a spatial biodiversity modeling pipeline, hereafter referred to as the SBM pipeline or the pipeline, for top-down modeling strategies that predict community-level biodiversity metrics. In addition, this Master's thesis also provides a case study demonstrating the applicability and flexibility of this pipeline. To achieve this objective, a pipeline is first developed that integrates generalized data preprocessing, feature engineering, model training, and evaluation steps for top-down models. This pipeline is designed to be flexible and adaptable to a wide range of datasets and modeling approaches. Next, a case study is conducted using this pipeline to compare the performance of different models in different contexts on derived biodiversity metrics from the United Kingdom Butterfly Monitoring Scheme (UKBMS) dataset \parencite{UKBMS2026}, accessed through the Global Biodiversity Information Facility (GBIF)\footnote{Global Biodiversity Information Facility (GBIF). Available at: \url{https://www.gbif.org/} (Accessed: 26 April 2026).}. Models used in this case study will include statistical models, specifically generalized linear models and generalized additive models \parencite{nelder1972,hastie1986}, as well as machine learning models, including random forests, gradient boosting, and artificial neural networks \parencite{lek1996, Breiman2001RF, friedman2001}. These results are then analyzed to provide insights into the strengths and weaknesses of different top-down models across different contexts. By applying the pipeline to the UKBMS dataset for rigorous model comparison, this thesis demonstrates the pipeline's flexibility and applicability to a range of community-level biodiversity modeling scenarios. This comparative framework also enables a systematic assessment of how different models (statistical or machine learning) perform under varying analytical conditions, thereby providing insight into their relative performance, strengths, and limitations. The key research questions that this Master's thesis aims to address are:

\begin{enumerate}
    \item How do different models perform compared to each other?
    \item How do models perform under different cross-validation strategies, specifically random cross-validation vs. spatial cross-validation?
    \item How do models perform in data-limited scenarios?
\end{enumerate} 

The remainder of this thesis is structured as follows. Section~\ref{sec:background} provides background on biodiversity, including how it is defined and measured in Section~\ref{subsec:biodiversity}, problem formulation of top-down modeling approaches in Section~\ref{subsec:modeling-definition}, a description of the modeling approaches used in this thesis in Section~\ref{subsec:Models-in-this-thesis}, and a literature review and connection to the research questions in Section~\ref{subsec:current-approaches}. Section~\ref{sec:Methods} describes the methodological framework, including the pipeline design in Section~\ref{subsec:pipeline}, the application of the pipeline to the UKBMS dataset in Section~\ref{subsec:application}, and how the models will be compared to one another in Section~\ref{subsec:model_comparison}. Section~\ref{sec:Results} presents the results of the case study and illustrates the flexibility of the proposed pipeline. Section~\ref{sec:discussion} discusses the implications of these results for biodiversity modeling and conservation, while Section~\ref{sec:ethics} addresses relevant ethical considerations. Section~\ref{sec:future_work} outlines potential future directions and applications for research in this area. Finally, Section~\ref{sec:conclusion} concludes the thesis by summarizing the key findings and their implications for biodiversity modeling and conservation.